\section{Introduction}
Procedural memory lets an agent reuse a successful method without deriving it again. Yet success is conditional. A record may have been complete before reaching a parser; a cache may have belonged to one request; a diagnostic object may have contained only public fields. If the target changes these assumptions, a memory can accurately describe the source task while providing an incomplete guide to the new one.

We ask: \emph{when a coding agent receives a source procedure in a target context where its assumptions no longer hold, how does that memory change functional completion and focal security failures?} A second question asks whether a generic applicability reminder changes this behavior. We supply a frozen correct memory, a matched irrelevant memory, correct memory with the reminder, or no memory. The target task and changed assumption remain fixed.

The original hypothesis was that correct source memory would increase functional completions that violate the target's focal security property. The completed experiment does not establish this as a general effect. Corrected estimates have both signs across six configurations. This mixed result becomes informative when completion, measurement availability, and implementation behavior are examined together. A lower failure rate can reflect an unfinished task. A witness pass can establish confidentiality without establishing useful output. Two implementations can fail the same witness through different mechanisms.

The study contributes (1) thirteen controlled synthetic families with explicit source assumptions and separate functionality and security observations; (2) corrected treatment comparisons, joint outcomes, and bounds that retain completed measurements within otherwise invalid runs; and (3) systematic exploratory coding of a transparently selected 52-run sample, complemented by matched transcript, patch, and outcome comparisons. The coding separates what agents say before editing from the protections they implement. The resulting contribution is a more precise account of observed memory transfer, including successful adaptations and counterexamples to harmful reuse.

All targets contain a context shift, and memories are supplied directly. The intervention therefore estimates the effect of supplied memory in those targets; it does not identify the effect of context shift or evaluate a retrieval system. Comparisons are made within each model and agent configuration.
